\section{Discrete symmetry breaking at low temperatures}

\begin{proposition}[The group $T_{C,\lambda}$ and Remark (18.22), Georgii (18.20)--(18.22)]
    \label{prop:dsb-18-22}
    \uses{def:lo-G0, def:lo-18-8}
    \lean{MeasureTheory.GibbsMeasure.latticeParity, MeasureTheory.GibbsMeasure.CubeTransformation, MeasureTheory.GibbsMeasure.CubeTransformation.onLattice, MeasureTheory.GibbsMeasure.CubeTransformation.onCube, MeasureTheory.GibbsMeasure.CubeTransformation.refl, MeasureTheory.GibbsMeasure.CubeTransformation.refl_onCube_toFun, MeasureTheory.GibbsMeasure.CubeTransformation.pureSpin, MeasureTheory.GibbsMeasure.CubeTransformation.planeRot, MeasureTheory.GibbsMeasure.CubeTransformation.planeRotSite, MeasureTheory.GibbsMeasure.CubeTransformation.planeRot_onCube_toFun_apply, MeasureTheory.GibbsMeasure.CubeTransformation.latticePattern_onLattice, MeasureTheory.GibbsMeasure.map_periodicGibbsField_periodize, MeasureTheory.GibbsMeasure.measurePreserving_of_mem_GZero, MeasureTheory.GibbsMeasure.CubeTransformation.IsSymmetry, MeasureTheory.GibbsMeasure.CubeTransformation.IsSymmetry.measurePreserving_onLattice_of_mem_GZero, MeasureTheory.GibbsMeasure.CubeTransformation.isSymmetry_refl, MeasureTheory.GibbsMeasure.measurePreserving_latticeRefl_of_mem_GZero, MeasureTheory.GibbsMeasure.latticePattern_shift_preimage_tauPow, MeasureTheory.GibbsMeasure.IsPhiSymmetric.const_mul}
    \leanok

    A cube transformation is a signed coordinate permutation of $\mathbb Z^d$ about the centre of
    the unit cube $C$ together with a $2$-periodic family of measurable bijections of $E$: the
    reflections $r_k$ of (17.14), the plane rotation $r_0$ of (18.20),
    $(r_0\omega)_i = \omega_{(i_2, 1-i_1, i_3, \dots)}$, and the pure spin transformations (18.21),
    $\tau_j = \tau_i$ for $j \equiv i \pmod 2$, are instances; the spatial part is more general
    than Georgii's group generated by $r_0, \dots, r_d$, and coincides with it for $d = 2$. It is
    a symmetry of the $C$-potential $\Phi$ when it fixes the cube interaction; $I_{C,\lambda}(\Phi)
    \subseteq I_{C,\lambda}(\beta\Phi)$. Remark (18.22): every $\mu \in \mathcal G_0(\Phi)$ is
    invariant under every symmetry, in particular under $r_1, \dots, r_d$, because the periodic
    Gibbs distributions on the torus are, and invariance passes to cluster points on local events
    (Example (5.20)(3)). No summability of $\Phi$ enters.
\end{proposition}
\begin{proof}
    \leanok
\end{proof}

\begin{proposition}[Stable partitions of the local ground states and the inclusion (18.24), Georgii (18.23)--(18.24)]
    \label{prop:dsb-18-24}
    \uses{prop:dsb-18-22, def:lo-18-6, lem:lo-ocean-tail}
    \lean{MeasureTheory.GibbsMeasure.faceRestrict, MeasureTheory.GibbsMeasure.IsPStablePartition, MeasureTheory.GibbsMeasure.IsPhiSymmetric, MeasureTheory.GibbsMeasure.IsPStablePartition.eq_of_faceRestrict_eq, MeasureTheory.GibbsMeasure.IsPStablePartition.pairwise_disjoint, MeasureTheory.GibbsMeasure.IsPStablePartition.eq_of_mem_latticePattern_add_single, MeasureTheory.GibbsMeasure.IsPStablePartition.eq_of_adj, MeasureTheory.GibbsMeasure.IsPStablePartition.exists_subset_latticePattern_of_connected, MeasureTheory.GibbsMeasure.latticeOceanPart_mono, MeasureTheory.GibbsMeasure.IsPStablePartition.exists_latticeOceanPart_nonempty, MeasureTheory.GibbsMeasure.IsPStablePartition.setOf_latticeOceanPart_nonempty_subset_iUnion, MeasureTheory.GibbsMeasure.IsPStablePartition.disjoint_setOf_latticeOceanPart_nonempty, MeasureTheory.GibbsMeasure.IsPStablePartition.latticeOceanPart_iUnion_eq, SimpleGraph.IsOceanIn.nonempty_inter, SimpleGraph.Connected.subset_image_val_supp_of_disjoint}
    \leanok

    A $P$-stable partition $\{G_1,\dots,G_N\}$ (18.23)(a) has pairwise disjoint projections of the
    $G_n$ on each $P$-face of $C$, and is $\Phi$-symmetric (18.23)(b) if any two members are
    exchanged by a symmetry. Adjacent sites of $V_P(G_\varepsilon(\Phi),\omega)$ share a $P$-face,
    so a connected subset of $V_P(G_\varepsilon,\omega)$ lies in a single $V_P(G_n,\omega)$, and an
    ocean of the former is an ocean of the latter: (18.24),
    $\{\xi^0_P(G_\varepsilon(\Phi),\cdot) \neq \emptyset\} \subseteq
    \bigcup_n\{\xi^0_P(G_n,\cdot) \neq \emptyset\}$, the events on the right being pairwise disjoint
    because two oceans in the plane meet (a disjoint one would lie in a single, finite, island of
    the other).
\end{proposition}
\begin{proof}
    \leanok
\end{proof}

\begin{theorem}[Georgii Theorem (18.25) and Comment (18.26)(1)]
    \label{thm:dsb-18-25}
    \uses{prop:dsb-18-24, thm:lo-18-17, lem:lo-ocean-tail, prop:lo-18-12}
    \lean{MeasureTheory.GibbsMeasure.preimage_shift_setOf_latticeOceanPart_nonempty_eq_tauPow, MeasureTheory.GibbsMeasure.preimage_onLattice_setOf_latticeOceanPart_nonempty, MeasureTheory.GibbsMeasure.measure_setOf_latticeOceanPart_nonempty_eq, MeasureTheory.GibbsMeasure.measure_setOf_latticeOceanPart_nonempty_pos, MeasureTheory.GibbsMeasure.le_cond_measure_setOf_zero_mem_latticeOceanPart, MeasureTheory.GibbsMeasure.mutuallySingular_cond_setOf_latticeOceanPart_nonempty, MeasureTheory.GibbsMeasure.isGibbsMeasure_cond_setOf_latticeOceanPart_nonempty, MeasureTheory.GibbsMeasure.measurePreserving_shift_cond_setOf_latticeOceanPart_nonempty, MeasureTheory.GibbsMeasure.measurePreserving_onLattice_cond_setOf_latticeOceanPart_nonempty, MeasureTheory.GibbsMeasure.map_onLattice_cond_setOf_latticeOceanPart_nonempty, MeasureTheory.GibbsMeasure.map_shift_cond_setOf_latticeOceanPart_nonempty, MeasureTheory.GibbsMeasure.exists_gibbsMeasure_family, MeasureTheory.GibbsMeasure.exists_gibbsMeasure_family_of_crossingBound_lt_one, MeasureTheory.GibbsMeasure.eventually_exists_gibbsMeasure_family, MeasureTheory.GibbsMeasure.le_encard_extremePoints_invariantG_shiftGroup_of_isRSymmetric, MeasureTheory.linearIndependent_of_measure_eq_one_of_measure_eq_zero, MeasureTheory.MeasurePreserving.map_cond_of_preimage_eq}
    \leanok

    In the plane, for a standard Borel state space, a confining cube interaction and a
    $\Phi$-symmetric $P$-stable partition of the local ground states, for $\beta$ large --- more
    precisely whenever $4z(t(G_\varepsilon(\Phi),\beta\Phi)) < 1$, which is Comment (18.26)(1) ---
    there are pairwise mutually singular, hence distinct, Gibbs measures
    $\mu_1,\dots,\mu_N \in \mathcal G(\beta\Phi)$ with (i) $\mu_n(\xi^0_P(G_n,\cdot) \neq \emptyset)
    = 1$ and $\mu_n(0 \in \xi^0_P(G_n,\cdot)) \geq 1 - 4z$, hence $\to 1$; (ii) $\mu_n$ invariant
    under the shifts $\theta_a$ with $r^aG_n = G_n$ and under the symmetries fixing $G_n$;
    (iii) $\tau(\mu_m) = \mu_n$ when $\tau G_m = G_n$, and $\theta_{e_k}(\mu_m) = \mu_n$ when
    $\tau = r_k$. Georgii's proof: some $\mu \in \mathcal G_0(\beta\Phi)$ has
    $\mu(\xi^0_P(G_\varepsilon) \neq \emptyset) \geq 1 - 4z > 0$ by (18.17), (18.24) and the symmetry
    give $\mu(A_1) = \dots = \mu(A_N) > 0$, the $A_n$ are tail events (Lemma
    \ref{lem:lo-ocean-tail}, for patterns that are not $r$-symmetric), so $\mu_n = \mu(\cdot\mid A_n)$
    is Gibbs by (7.7)(b), and invariance passes to the conditional measures. The count
    $N \leq |\operatorname{ex}\mathcal G_{\tilde\Theta}(\beta\Phi)|$ is proved for
    $\tilde\Theta = \Theta$, i.e. when all $G_n$ are $r$-symmetric, from (14.18) and the linear
    independence of measures separated by disjoint events; for a proper finite-index subgroup
    $\tilde\Theta$ it needs the $(\mathcal P_{\tilde\Theta},\mathcal I)$-kernel of the remark below
    (14.18), not in the tree. $N > 1$ is not required.
\end{theorem}
\begin{proof}
    \leanok
\end{proof}
